\section{Автоматическая векторизация программного кода}
\label{sec:Chapter2} \index{Chapter2}

\subsection{SIMD как аппаратная основа векторизации}

SIMD представляет собой модель вычислений, при которой одна инструкция одновременно применяется к нескольким элементам данных, размещённым в специальном векторном регистре. Такой подход позволяет существенно увеличить вычислительную производительность при выполнении однотипных операций над независимыми данными.

Принцип работы SIMD можно проиллюстрировать на примере сложения двух массивов целых чисел. При скалярном выполнении операция сложения производится последовательно для каждой пары элементов:
\[
c_i = a_i + b_i,\quad i = 0,\dots,n-1
\]
При использовании SIMD несколько элементов обрабатываются одновременно одной инструкцией:
\[
c_{[i,i+m]} = a_{[i,i+m]} + b_{[i,i+m]},\quad i = 0,\dots,n/m-1,
\]
где \(m\) — количество данных, помещающихся в векторный регистр. Если ширина векторного регистра составляет 128 бит, то за одну инструкцию можно выполнить сложение двух 64-битных целых чисел (\(m=2\)) или четырех 32-битных (\(m=4\)).

Современные процессорные архитектуры предоставляют развитую поддержку SIMD-вычислений. В архитектуре ARM данная функциональность предоставляется расширениями NEON и SVE, отличающимися шириной векторных регистров и набором поддерживаемых операций. В архитектуре x86 аналогичные возможности реализованы через наборы инструкций SSE, AVX, AVX2 и AVX-512.

Использование SIMD особенно эффективно в задачах:
\begin{itemize}
    \item обработки массивов числовых данных;
    \item цифровой обработки сигналов;
    \item криптографических вычислениях;
    \item мультимедийной обработки;
    \item математических вычислений.
\end{itemize}

Однако эффективное применение SIMD-инструкций предъявляет ряд требований к структуре программы. Операции должны обладать независимостью по данным, иметь совместимый тип и размер операндов, а также допускать группировку в векторные операции.

Кроме того, непосредственное использование SIMD-инструкций связано с рядом практических трудностей. Разработчику необходимо учитывать особенности конкретной аппаратной архитектуры, различия в наборах инструкций и специфику работы с векторными регистрами. Это существенно усложняет разработку и снижает переносимость программного обеспечения.

По этой причине особую важность приобретают механизмы автоматической векторизации, позволяющие компилятору самостоятельно выявлять участки кода, пригодные для преобразования в векторную форму.

\subsection{Подходы к автоматической векторизации}

Автоматическая векторизация представляет собой класс оптимизаций компилятора, направленных на преобразование скалярного программного кода в эквивалентную форму, использующую SIMD-инструкции.

В современных компиляторах можно выделить два основных подхода к автоматической векторизации:
\begin{itemize}
    \item цикловая векторизация;
    \item SLP-векторизация.
\end{itemize}

\subsubsection{Цикловая векторизация}

Цикловая векторизация ориентирована на поиск параллелизма внутри циклов. Рассмотрим следующий пример:
\[
\text{for } i := 0; i < n; i++ \{\; 
    c[i] = a[i] + b[i]
\}
\]
Такой цикл содержит естественный параллелизм по итерациям. Компилятор может преобразовать его таким образом, чтобы за одну итерацию выполнялась обработка сразу нескольких элементов.

Эффективность цикловой векторизации во многом определяется возможностью доказать отсутствие зависимостей между итерациями цикла. Для этого компилятору необходимо выполнять сложный статический анализ, включающий:
\begin{itemize}
    \item анализ псевдонимов памяти (alias analysis);
    \item анализ межитерационных зависимостей (loop-carried dependency analysis);
\end{itemize}
а также ряд вспомогательных преобразований, таких как:
\begin{itemize}
    \item разворачивание цикла (loop unrolling);
    \item выделение граничных итераций (loop peeling).
\end{itemize}
Реализация подобных механизмов требует развитой инфраструктуры анализа и существенно усложняет оптимизационный конвейер компилятора.

\subsubsection{SLP-векторизация}

SLP-векторизация ориентирована на поиск параллелизма внутри ациклических участков программы. В отличие от цикловой векторизации, данный подход не анализирует зависимости между итерациями, а выявляет независимые однотипные скалярные операции, которые могут быть объединены в векторную инструкцию.

Например, последовательность:
\[
x[0] = a[0] + b[0],\quad x[1] = a[1] + b[1]
\]
может быть преобразована в одну векторную операцию сложения.

SLP-векторизация особенно эффективна в случаях, когда параллелизм явно присутствует в структуре вычислений, но не выражен через циклические конструкции.

\subsubsection{Сравнение подходов}

Цикловая и SLP-векторизация существенно различаются по сложности реализации. Цикловая векторизация требует сложного анализа, но позволяет оптимизировать циклические участки кода, которые зачастую являются основной вычислительной нагрузкой большого класса приложений. SLP-векторизация обладает меньшей вычислительной сложностью и проще интегрируется в инфраструктуру компилятора.

Для компилятора Go, где отсутствуют циклические оптимизации и нужная для них инфраструктура анализа, а также где одним из ключевых требований является быстрое время компиляции, SLP-подход представляет особый практический интерес.

\subsection{Метод Superword-Level Parallelism}

Метод Superword-Level Parallelism был предложен в работе Samuel Larsen и Saman Amarasinghe~\cite{Larsen2004} как способ выявления параллелизма на уровне независимых скалярных операций внутри базовых блоков программы. Основная идея метода заключается в объединении нескольких однотипных независимых скалярных операций в одну векторную.

\subsubsection{Основные этапы алгоритма}

Классический алгоритм SLP-векторизации~\cite{Larsen2004} состоит из нескольких последовательных этапов, выполняемых в пределах одного базового блока.

\begin{enumerate}
    \item \textbf{Поиск начальных пар}. Алгоритм сканирует инструкции базового блока и ищет пары соседних обращений к памяти – смежные загрузки или сохранения данных. Такие операции могут послужить естественным началом и/или завершением параллельных последовательностей вычислений.
    
    \item \textbf{Расширение групп}. От найденных начальных пар алгоритм итеративно строит новые группы, анализируя поток данных:
    \begin{itemize}
        \item \textit{Расширение вверх}: для каждой пары инструкций в группе анализируются их аргументы – если находятся две однотипные независимые операции, они могут образовать новую пару.
        \item \textit{Расширение вниз}: рассматриваются операции, использующие результат исходных операций. Среди нескольких возможных вариантов выбирается та пара использований, которая приносит наибольшую выгоду с точки зрения модели стоимости.
    \end{itemize}
    Процесс продолжается, пока находятся новые пары, удовлетворяющие условиям независимости и изоморфности операций.
    
    \item \textbf{Объединение пар в группы}. После завершения расширения ищутся пары, которые можно объединить. Если в двух парах есть общая операция, то их можно объединить в тройки, четвёрки и т.д., формируя потенциальные векторные регистры большей ширины.
    
    \item \textbf{Планирование}. Векторные инструкции должны располагаться в корректном порядке с учётом зависимостей по данным. Выполняется топологическая сортировка групп: ни одна группа не может быть помещена в выходной поток раньше тех групп, от которых она зависит через аргументы.
    
    \item \textbf{Генерация векторного кода}. Для каждой группы формируются соответствующие векторные инструкции целевой архитектуры. При необходимости группы разбиваются до поддерживаемой аппаратурой ширины. Если какая-то часть промежуточного результата используется вне векторизованного фрагмента, то нужное значение скаляризуется – вставляются специальные операции извлечения отдельных данных из векторного регистра в обычный.
\end{enumerate}

Такой набор этапов гарантирует корректность векторизации и позволяет последовательно расширять функциональность, начиная с простых парных операций.

\subsubsection{Преимущества метода SLP}

Метод Superword-Level Parallelism обладает рядом существенных преимуществ.
\begin{itemize}
    \item \textbf{Простота и низкие накладные расходы}. Алгоритм работает исключительно на уровне отдельных базовых блоков, что упрощает реализацию и снижает время компиляции.
    \item \textbf{Расширяемость}. Метод SLP допускает поэтапное наращивание функциональности: от поддержки простейших целочисленных операций до межблочной векторизации. Более того, SLP может быть расширен в сторону учёта циклических конструкций. Примером служит Loop-Aware SLP, реализованный в GCC~\cite{loopAwareSLP}, который анализирует возможности векторизации инструкций через границы итераций цикла, объединяя преимущества классического SLP и цикловой векторизации.
\end{itemize}
Перечисленные свойства делают SLP прагматичным компромиссом между сложностью, временем компиляции и достигаемым ускорением.

\subsubsection{Ограничения метода}

Несмотря на преимущества, SLP обладает рядом ограничений.
\begin{itemize}
    \item Метод эффективен только при наличии явно выраженных групп независимых однотипных операций. Если вычисления имеют сложную структуру зависимостей или требуют существенного переупорядочивания данных, эффективность SLP снижается.
    \item Алгоритм чувствителен к качеству модели стоимости. Слишком агрессивная векторизация может привести к увеличению числа вспомогательных операций упаковки/распаковки и ухудшению производительности.
\end{itemize}

\subsubsection{Причины выбора SLP для компилятора Go}

Выбор метода SLP для реализации в компиляторе Go обусловлен рядом факторов. Данный подход хорошо согласуется с существующей архитектурой компилятора. Реализация SLP требует существенно меньшего объёма инфраструктурных изменений по сравнению с цикловой векторизацией. Алгоритм обладает умеренной вычислительной сложностью, что соответствует требованию сохранения малого времени компиляции.

Таким образом, SLP представляет собой практически обоснованный компромисс между сложностью реализации и потенциальным приростом производительности.

\subsection{Выводы по главе}
В данной главе были рассмотрены аппаратные основы SIMD-вычислений, основные подходы к автоматической векторизации и детали метода Superword-Level Parallelism.
Показано, что SLP-векторизация, в отличие от цикловой, не требует анализа итерационных зависимостей и сложных трансформаций циклов. Это делает её практически применимым подходом для интеграции в компилятор Go: умеренная алгоритмическая сложность сочетается с возможностью ускорять ацикличный код, типичный для многих серверных и системных программ на Go.
